\subsection{相对于子空间的连续性}

设 X 与 Y 是两个拓扑空间，f 是 X 到 Y 中的一个映射，B 是 Y 的包含 \(f(\mathbf{X})\) 的一个子集。诱导拓扑作为初始拓扑的定义（§ 2，第 3 目，命题 4）表明，f 在 \(x \in X\) 处连续，当且仅当 X 到 Y 的子空间 B 中的、与 f 有同一图象的那个映射在 x 处连续。

现在设 A 是 X 的一个子集；若 f 在 \(x \in A\) 处连续（相应地，在 X 上连续），则它的限制 \(f|A\) 是子空间 A 到 Y 中的一个映射，由 §2 第 1 目命题 2，它在 x 处连续（相应地，在 A 上连续）。我们有时会说，映射 \(f: X \to Y\) 在 \(x \in A\) 处相对于 A 连续（相应地，相对于 A 连续），如果它的限制 \(f|A\) 在 x 处连续（相应地，在 A 上连续）。

应当注意，\(f|A\) 可以连续而 \(f\) 不在 \(X\) 的任何点连续；这种现象的一个例子由特征函数 \(\varphi_A\)（子集 \(A\)（\(X\) 的子集）的特征函数，该子集使 \(A\) 与它的余集都在 \(X\) 中稠密（§ 2，习题 11））给出，其中 \(\varphi_A\) 被看作 \(X\) 到离散空间 \(\{0, 1\}\) 中的一个映射。\(\varphi_A\) 不在 \(X\) 的任何点连续，但它在 \(A\) 上的限制是常数，因而是连续的。

若 A 是点 \(x \in A\) 在 X 中的一个邻域，且 \(f: X \to Y\) 使 \(f|A\) 在 x 处连续，则 f 在 x 处连续；因为 x 在 A 中的每个邻域都是 x 在 X 中的邻域（连续性的局部特征）。

命题 4. 设 \((\mathbf{A}_t)_{i\in I}\) 是拓扑空间 \(\mathbf{X}\) 的子集的一个族，其诸内部覆盖 \(\mathbf{X}\)，或者它是 \(\mathbf{X}\) 的一个局部有限闭覆盖。设 \(f\) 是 \(\mathbf{X}\) 到拓扑空间 \(\mathbf{X}'\) 中的一个映射。若 \(f\) 在每个子空间 \(\mathbf{A}_t\) 上的限制都连续，则 \(f\) 连续。

因为，若 \(\mathbf{F}'\) 是 \(\mathbf{X}'\) 的一个闭子集且 \(\mathbf{F} = \overline{f}^{1} (\mathbf{F}')\)，则 \(\mathbf{F} \cap \mathbf{A}_i\) 在 \(\mathbf{A}_i\) 中是闭集（对每个 \(i \in I\)）（\(\S 2\)，第 1 目，定理 1），因而由第 1 目命题 3，\(\mathbf{F}\) 在 \(\mathbf{X}\) 中是闭集；于是结论由 \(\S 2\) 第 1 目定理 1 得出。

